% Generated draft: requires mathematical and source review.
\documentclass[11pt]{article}
\usepackage[margin=25mm]{geometry}
\usepackage{fontspec}
\setmainfont{Latin Modern Roman}
\usepackage{amsmath,amssymb,mathtools,unicode-math}
\setmathfont{Latin Modern Math}
\usepackage{enumitem,needspace,xurl,hyperref}
\setlength{\parindent}{0pt}
\setlength{\parskip}{5pt}
\setlength{\emergencystretch}{3em}
\newcommand{\N}{\mathbb{N}}
\newcommand{\Z}{\mathbb{Z}}
\newcommand{\Q}{\mathbb{Q}}
\newcommand{\R}{\mathbb{R}}
\newcommand{\C}{\mathbb{C}}
\newcommand{\F}{\mathbb{F}}
\newcommand{\A}{\mathbb{A}}
\newcommand{\PP}{\mathbb{P}}
\newcommand{\E}{\mathbb{E}}
\newcommand{\eps}{\varepsilon}
\newcommand{\dd}{\,\mathrm{d}}
\begin{document}
\section*{UnsolvedMath ID: 701163}
\textbf{Source problem code:} catalogue entry
\subsection{Definitions and mathematical statement}
Let \(F(x,y)\) be the free group on two generators. For this draft, use the commutator convention \([a,b]=a^{-1}b^{-1}ab\). The supplied record does not specify its commutator convention; compatibility with the original source requires review. Define words \(u_n\in F(x,y)\), for integers \(n\geq 1\), recursively by
\[
u_1(x,y)=[x,y],\qquad
u_{n+1}(x,y)=[[u_n(x,y),x],[u_n(x,y),y]].
\]
For a group \(G\), the notation \(G\models u_n=1\) means
\[
\forall a,b\in G,\qquad u_n(a,b)=1_G.
\]
Thus the index \(n\) in such an identity is uniform over all pairs of elements of \(G\). Define the derived series by \(G^{(0)}=G\) and
\[
G^{(r+1)}=\langle [a,b]:a,b\in G^{(r)}\rangle
\quad(r\geq 0).
\]
The group \(G\) is soluble if \(G^{(r)}=\{1_G\}\) for some integer \(r\geq 0\).

\textbf{Part (b), the remaining question in the supplied record.} Is it true that, for every finite group \(G\),
\[
G\text{ is soluble}
\quad\Longleftrightarrow\quad
\exists n\geq 1\ \forall a,b\in G,\quad u_n(a,b)=1_G?
\]

\textbf{Part (a), included as archived solved context.} Does there exist a sequence of two-variable group words \(w_1,w_2,\ldots\in F(x,y)\) such that each identity \(w_n=1\) implies \(w_{n+1}=1\), and, for every finite group \(G\),
\[
G\text{ is soluble}
\quad\Longleftrightarrow\quad
\exists n\geq 1\ \forall a,b\in G,\quad w_n(a,b)=1_G?
\]
Here \(w_n\) renames the words called \(u_n\) in part (a), to distinguish this existentially chosen sequence from the prescribed recursion in part (b). The source does not explicitly specify the class of groups over which implication of successive identities is required. Its two evident readings are
\[
\forall H\in\mathcal C\ \forall n\geq 1,\qquad
(H\models w_n=1)\Longrightarrow(H\models w_{n+1}=1),
\]
with \(\mathcal C\) either all groups or all finite groups; this scope remains unresolved here. The supplied archive reports an affirmative answer to part (a). This report does not establish the assertion in part (b).
\subsection{Short English statement}
Problem 15.75(b), proposed by B. I. Plotkin, asks whether a finite group is soluble exactly when some word in the specified recursive sequence vanishes on every pair of its elements. The index may depend on the group. The supplied record marks the overall problem partially solved and reports that part (a), concerning the existence of a suitable sequence, has been solved affirmatively. These are supplied source assessments, not independently verified conclusions.
\subsection{Sources}
\begin{itemize}
\item Local UnsolvedMath registry: supplied complete source record, permanent ID \(701163\), exact source code \texttt{catalogue entry}, title \emph{Kourovka Notebook Problem 15.75}. Upstream revision: \texttt{d371ee917ac9647fbb06cd6fc529deec6f3dda82}.
\item B. I. Plotkin, proposer of Problem 15.75, \emph{The Kourovka Notebook (Unsolved Problems in Group Theory)}, 15th issue (2002); the supplied background identifies an update in the 21st issue (2026), \texttt{arXiv:1401.0300v48}. The supplied collection description names E. I. Khukhro and V. D. Mazurov as editors.
\item The supplied Archive of Solved Problems entry reports that the sequences requested in part (a) exist, citing T. Bandman, F. Grunewald, G.-M. Greuel, B. Kunyavskii, G. Pfister, and E. Plotkin, \emph{Compositio Math.} \textbf{142} (2006), 734--764.
\item The same archive entry cites J. N. Bray, J. S. Wilson, and R. A. Wilson, \emph{Bull. Lond. Math. Soc.} \textbf{37} (2005), 179--186.
\item The same archive entry cites E. Ribnere, \emph{Monatsh. Math.} \textbf{157} (2009), 387--401. The author spelling is retained as supplied.
\item The supplied status review is dated \texttt{2026-10-16}, labels the record \texttt{partially\_solved}, and states that solved parts appear in the archive while remaining parts appear in the main section. This is supplied dated curation, not an independently verified literature assessment. No external source was visited for this draft.
\end{itemize}
\end{document}
